\begin{table}[htbp]
\centering
\caption{\textbf{Unit of replication in the perturbation analysis.} Transcription factors with the largest cell-level response.}
\label{tab:cell_level}
\smallskip
\small
\begin{tabular}{lrrrrrc}
\toprule
Target & Cells & Responding (cell) & Responding (position) & Inflation & ICC & Self detected \\
\midrule
GATA1 & 95 & 30 & 28 & 0.9 & 0.000 & no \\
MAX & 161 & 7 & 4 & 0.6 & 0.000 & yes \\
MYBL2 & 81 & 1 & 0 & 0.0 & 0.000 & no \\
TBP & 143 & 0 & 0 & 0.0 & 0.000 & no \\
SRF & 120 & 0 & 0 & 0.0 & 0.000 & no \\
THAP1 & 107 & 0 & 0 & 0.0 & 0.000 & no \\
TFDP1 & 140 & 0 & 0 & 0.0 & 0.000 & no \\
CEBPZ & 104 & 0 & 0 & 0.0 & 0.000 & no \\
STAT5A & 170 & 0 & 0 & 0.0 & 0.000 & no \\
E2F6 & 140 & 0 & 0 & 0.0 & 0.000 & no \\
THAP11 & 60 & 0 & 0 & 0.0 & 0.000 & no \\
TERF1 & 72 & 0 & 0 & 0.0 & 0.000 & no \\
BRF2 & 170 & 0 & 0 & 0.0 & 0.000 & no \\
ATF4 & 59 & 0 & 0 & 0.0 & 0.000 & no \\
SP2 & 115 & 0 & 0 & 0.0 & 0.000 & no \\
SETDB1 & 123 & 0 & 0 & 0.0 & 0.000 & no \\
HINFP & 157 & 0 & 0 & 0.0 & 0.000 & no \\
BRCA1 & 87 & 0 & 0 & 0.0 & 0.000 & no \\
DNMT1 & 161 & 0 & 0 & 0.0 & 0.000 & no \\
ERCC2 & 116 & 0 & 0 & 0.0 & 0.000 & no \\
\bottomrule
\end{tabular}
\\[2pt]
\footnotesize Responding features under a cell-level Mann--Whitney test against a position-pooled test on exactly the same cells. The intraclass correlation is the between-cell share of feature-activation variance (median across factors 0.000). Median inflation of the responding-feature count under position pooling: 0.0$\times$. The last column records whether the knocked-down gene itself appears in the predicted set, a positive control for detection power.
\end{table}
